\section{A1 Reuse total decoding for raw parent codes}
\label{sec:parents}

\textbf{Target.} \src{AnnotatedParents.lean}{49}{97},
\code{check_raw_iff}. The theorem connects a numeric parent-list checker to
an existential list of indices. Its forward proof reconstructs that list
through a 44-line proof-body induction, despite the preceding
\src{AnnotatedParents.lean}{29}{45}, \code{check_lists}, already establishing
the exact typed-list characterization.

\textbf{Change.} Choose an encoded list using the existing raw-table decoder,
rewrite to its exact recoding, then apply \code{check_lists}. The decoder
returns arbitrary raw six-field rows; their row codes provide the needed
natural-number list. This does not assert that the rows are valid proof lines.

\begin{Lean}
  constructor
  · intro h
    obtain ⟨rows, hr⟩ :=
      AnnotatedRenameScan.decodeTable_total parents parents
        (Nat.le_refl _)
    have he : AnnotatedCodeArithmetic.listCode
        (rows.map AnnotatedRename.rowCode) = parents :=
      AnnotatedRenameScan.tableCode_of_decode hr
    refine ⟨rows.map AnnotatedRename.rowCode, he.symm, ?_⟩
    exact (check_lists fuel history expected _ hf).mp
      (by simpa only [he] using h)
  · rintro ⟨indices, rfl, h⟩
    exact (check_lists fuel history expected indices hf).mpr h
\end{Lean}

This is the candidate proof body, line-wrapped for the page. The full
statement and non-conflicting \code{check_raw_iff_short} declaration are in
\pathcode{annotated/AnnotatedCandidates.lean}. The exact conservative patch
is \pathcode{annotated/proposed.patch}; it has not been applied to the
supplied source snapshot.

\subsection{Why the bridge is available}

\src{AnnotatedRenameScan.lean}{124}{162} proves both
\code{tableCode_of_decode} and \code{decodeTable_total}. Supplying
\code{parents} as both code and sufficient external-natural fuel yields
\code{rows}; exact recoding gives \code{he}. The parent-checker's forward
direction is now the existing typed-list equivalence applied after that
rewrite. The backward direction already has the encoded list, so it needs
no decoder or reconstruction at all.

The import chain is already present. \code{AnnotatedParents} reaches
\code{AnnotatedRenameScan} through \code{AnnotatedSchemaMaps} and
\code{AnnotatedRenameScanPR}. There is no new import and no need for a
new list-code injectivity theorem.

\subsection{What must not be conflated}

The decoder is used in a \code{Prop} proof; it is not inserted into the
executable checker or its primitive-recursive program. The hypothesis
\code{hf : expected.length ≤ fuel} and the checker fuel remain unchanged.
The original source's emphasis on avoiding an unbounded decoder is not a
reason to claim this dependency is free: the candidate deliberately uses
an existing \emph{bounded-decoder totality theorem} at a sufficient
external-natural fuel. If the intended architecture forbids that proof
dependency, keep the original induction even if this candidate elaborates.

\textbf{Static outcome and gate.} The whole declaration changes from 49 to
15 nonblank lines, and from 2,083 to 829 UTF-8 bytes, under the artifact's
counting rule. First check implicit arguments and the proof-local
\code{he} rewrite in the pinned Lean environment; then compare the old and
new axiom dependency lists. No kernel-term or runtime reduction is claimed.
